The use and handling of the various data has a direct impact on the \gls{DNP} group parameters.
For example, Huynh et al. compare the effects of cumulative fission yield and emission probabilities from the \gls{JEFF}, \gls{ENDF}, and \gls{JENDL} on the total delayed neutron yield \cite{huynh2014calculation}.
Mathieu et al. investigate different datasets as well, including the \gls{NNDC} website and various data compilations \cite{mathieu2012new, pfeiffer2002status, audi2003nubase, rudstam1993delayed}.
This can be expanded upon by comparing the \gls{IAEA} database, \gls{NNDC} website, \gls{JEFF}, \gls{ENDF}, \gls{JENDL}, \gls{ENSDF}, and various data compilations for various combinations of emission probabilities, half-lives, fission yields, and cross sections.

A relevant analysis to be conducted is comparing the cumulative fission yield to the full depletion approach.
Previous microscopic \gls{DNP} group parameter evaluations have used cumulative fission yields rather than depletion simulations.
This approach does not capture the full time dependence of decay chains leading to \glspl{DNP}, which could result in discrepancies when including the \gls{MSR} physical phenomena.
However, because using the cumulative fission yields would remove the need for depletion simulations, it is useful to determine whether the cumulative fission yield approach can be used.
It is also useful because the spatially resolved method is much more efficient using the cumulative fission yield rather than full spatially resolved depletion.

Another analysis of interest in regards to the data is on the number of \glspl{DNP} needed in the spatially resolved method to generate accurate results.
The computational cost and memory requirements for the spatially resolved method scale with the number of nuclides included, so minimizing them is of interest.
The impact of each \gls{DNP} is given in Equation \eqref{eq:delnu-net-nuc}, replicated here:

\begin{equation}
    \bar{\nu}_{d, i} = P_{n, i} \lambda_i N_{i, 0}.
    \nonumber
\end{equation}
Dividing $\bar{\nu}_{d, i}$ by the total delayed neutron yield, $\bar{\nu}_d$, provides the relative contribution of each \gls{DNP} to the total delayed neutron yield.
The \glspl{DNP} with the most impact are modeled preferentially to determine the minimum number of \glspl{DNP} to model with negligible changes to the \gls{DNP} group parameters.

\begin{comment}
Finally, the pulse and saturation irradiation \gls{DNP} group parameters can be combined in a different manner.
The Levenberg-Marquardt method is discussed in the literature for combining the pulse and saturation irradiation data \cite{levenberg1944method, marquardt1963algorithm, mills2014calculation}.
This method of fitting \gls{DNP} group parameters to both sets of data simultaneously may provide more accurate results.
\end{comment}